\chapter[Orientation and scope]{Generalized hierarchical Bayesian segmentation: problem, outputs, and scope}
\label{ch:orientation}
\chapterpromise{The research problem, the meaning of generality and hierarchy, the returned posterior quantities, and the conditions for exact computation.}

\begin{objectives}
\begin{itemize}
\item State the generalized hierarchical Bayesian segmentation problem and its principal outputs.
\item Explain why the partition recursions are common across observation models once segment marginal likelihoods are available.
\item Identify the roles of irregular designs, multiple sequences, known groups, and latent groups.
\item Distinguish exact conjugate calculations from explicitly approximate nonconjugate calculations.
\end{itemize}
\end{objectives}

Multiple-changepoint analysis partitions an ordered set of observations into contiguous segments
within which a statistical parameter is constant or follows a specified segment model. Given
responses $y_1,\ldots,y_n$ observed at ordered coordinates $u_1<\cdots<u_n$, the inferential
questions include the number of segments, the changepoint locations, the segment-specific
parameters, and uncertainty in all of these quantities. Bayesian formulations place a prior on the
partition and segment parameters and therefore yield posterior probabilities for competing
segmentations rather than only a single penalized estimate \citep{barry1992ppm,barry1993bayesCP,
fearnhead2006exact,hutter2006bpcr}.

The central statistical structure is a product-partition model. Conditional on a partition,
observations in different segments are independent after segment parameters are introduced. When a
segment parameter can be integrated out, each candidate segment $(i,j]$ contributes a segment
marginal likelihood. A prior that factorizes over segments or transitions then permits exact finite
summation over all admissible ordered partitions by dynamic programming. This yields marginal
likelihoods, posterior distributions for the segment count, ordered-changepoint marginals,
changepoint-event probabilities, and posterior moments of segment functionals
\citep{barry1992ppm,fearnhead2006exact,hutter2006bpcr}.

\BayesBreak develops this calculation as a generalized hierarchical Bayesian segmentation method.
The term \emph{generalized} refers to the following mathematical property: the recursion over
partitions does not depend on the particular observation family once the required segment marginal
likelihoods and posterior moments have been computed. For a regular exponential family with a
proper conjugate prior and finite normalizing constants, these segment quantities are available by
exact integration from sufficient statistics. Gaussian observations with an unknown segment mean,
Poisson counts with exposure, and Binomial observations with known trial counts are examples. A
nonconjugate segment model may use numerical integration or another stated approximation, but exact
partition summation for an approximate score array is not the same as exact Bayesian inference for
the intended observation model.

The term \emph{hierarchical} refers to probability models that couple sequences or groups through
shared partition structure or group-level parameters. The book treats: (i) a single ordered
sequence; (ii) irregularly spaced observations through design-dependent segment cohesion and
boundary hazards; (iii) multiple aligned sequences with common changepoints and conditionally
independent segment likelihoods; (iv) known groups of sequences; and (v) a finite latent-group
criterion in which sequences are assigned to candidate changepoint templates. The first four are
Bayesian models with explicit joint distributions under their stated assumptions. The finite
latent-group extension is presented according to its actual optimization criterion and is not
relabelled as a normalized mixture likelihood.

\paragraph{Principal contributions.}
\begin{itemize}
\item \textbf{Exponential-family segment integration.} A common notation separates sufficient
statistics, exposure or trial descriptors, and optional likelihood powers. Proper conjugate priors
give exact segment marginal likelihoods and posterior moments through ratios of normalizing
constants.

\item \textbf{Posterior computation over ordered partitions.} Forward and backward sum-product
recursions compute fixed-count marginal likelihoods, $p(k\mid y)$, ordered-changepoint marginals,
changepoint-event probabilities, and model-averaged segment functionals. A distinct max-sum
recursion with backtracking computes the joint MAP partition.

\item \textbf{Irregular-design priors.} Segment cohesion and interior-boundary hazard are separated
mathematically. This permits priors based on physical segment length or local design spacing without
confusing design geometry with exposure in the observation likelihood.

\item \textbf{Multiple-sequence hierarchies.} For aligned conditionally independent sequences with
a common partition, sequence-specific segment marginal likelihoods multiply. Known groups use the
same calculation within each group. A finite-candidate consistency result states conditions under
which the posterior concentrates on the common partition as the number of informative sequences
increases.

\item \textbf{Grouped and latent-group segmentation.} Known group memberships lead to groupwise
shared-changepoint posteriors. For unknown memberships, a finite latent-group objective alternates
auxiliary allocation weights with exact responsibility-weighted max-sum updates of the group
templates; monotonicity follows from a Jensen minorization argument.

\item \textbf{Numerical segment integration and prediction.} For nonconjugate observation models,
the effect of a uniform segmentwise log marginal-likelihood error is propagated through the
partition sums. Posterior prediction is defined by the fitted observation family and coordinate
support, and changepoint metrics use one-to-one matching.
\end{itemize}

\section{Research question and returned quantities}
The research question is whether one generalized Bayesian formulation can combine
observation-model-specific segment integration with posterior computation over contiguous
partitions while retaining irregular-design and hierarchical multiple-sequence structure. For a
sequence $y$, the principal returned quantities are
\[
\left\{p(y\mid k),\ p(k\mid y),\ p(t_p=h\mid y,k),\
p(b_h=1\mid y),\ \mathbb E[m(\theta_s)\mid y],\
\widehat t^{\mathrm{MAP}}(k)\right\}.
\]
These quantities have different meanings. In particular, a marginal mode of each changepoint
coordinate need not form the joint MAP partition, and a posterior mean curve is not a selected
piecewise-constant partition.

\begin{keyidea}
The observation model determines the segment marginal likelihoods and posterior moments. The
partition model determines how these quantities are summed over ordered segmentations. Generality
comes from separating these two mathematical calculations while keeping their probability models
explicit.
\end{keyidea}

\begin{figure}[H]
\centering
\resizebox{0.96\textwidth}{!}{\begin{tikzpicture}[node distance=8mm and 10mm]
\node[primarybox,text width=31mm] (data) {Ordered observations\\$u,y,d,a$};
\node[secondarybox,right=of data,text width=35mm] (segment) {Observation-model calculation\\$\log A^{(r)}_{ij}$};
\node[primarybox,above right=7mm and 13mm of segment,text width=32mm] (sum) {Sum-product recursion\\marginal likelihoods and posterior marginals};
\node[primarybox,below right=7mm and 13mm of segment,text width=32mm] (max) {Max-sum recursion\\joint MAP partition};
\node[successbox,right=13mm of sum,text width=31mm] (post) {$p(k\mid y)$, changepoint probabilities, Bayes curves};
\node[successbox,right=13mm of max,text width=31mm] (map) {Joint MAP changepoints and fitted segments};
\draw[arrPrimary] (data)--(segment);
\draw[arrPrimary] (segment)--(sum);
\draw[arrPrimary] (segment)--(max);
\draw[arrPrimary] (sum)--(post);
\draw[arrPrimary] (max)--(map);
\node[draw=Rule,dashed,rounded corners,fit=(sum)(max),inner sep=4mm,label=above:{\small dynamic programming over ordered partitions}] {};
\end{tikzpicture}
}
\caption{Observation-model calculations and inference over ordered partitions. Sum-product and
max-sum recursions use the same prior-adjusted segment scores but return different posterior or
optimization quantities.}
\label{fig:block-to-global}
\end{figure}

\section{Conditions and exclusions}
The exact results concern offline segmentation on a finite ordered candidate set with a prior and
likelihood that factor into segment or transition contributions. The method is not an online
run-length filter, a general nonparametric detector, or a proof that arbitrary dependence can be
represented without enlarging the state space. Exact Bayesian inference requires exact segment
marginal likelihoods. When these quantities are approximated, the resulting partition calculation
is exact only for the supplied approximate array; comparison with the intended posterior requires
additional numerical error analysis.

\begin{figure}[H]
\centering
\resizebox{0.96\textwidth}{!}{\begin{tikzpicture}[node distance=8mm and 10mm]
\node[successbox,text width=29mm] (exact) {Exact conjugate segment marginal likelihoods};
\node[secondarybox,below=of exact,text width=29mm] (approx) {Numerically approximated segment marginal likelihoods};
\node[primarybox,right=of exact,yshift=-8mm,text width=30mm] (dp) {Posterior and MAP calculations over ordered partitions};
\node[successbox,right=of dp,yshift=13mm,text width=28mm] (real) {Executed results with valid interpretation};
\node[neutralbox,right=of dp,yshift=-13mm,text width=28mm] (planned) {Planned validation studies};
\node[draw=Warning,fill=WarningPale,rounded corners,right=of dp,text width=30mm,align=center,minimum height=10mm] (excluded) {Executed results excluded from a stated analysis};
\draw[arrPrimary] (exact)--(dp);
\draw[arrSecondary,dashed] (approx)--node[above,sloped,font=\scriptsize]{with stated numerical assumptions}(dp);
\draw[arrPrimary] (dp)--(real);
\draw[arr,dashed] (dp)--(planned);
\draw[arr,draw=Warning] (dp)--(excluded);
\node[below=2mm of excluded,font=\scriptsize,text=Warning,align=center,text width=34mm] {The output remains recorded but does not support the incompatible claim.};
\end{tikzpicture}
}
\caption{Distinction among exact calculations, numerical approximations, executed empirical
results, planned validation, and executed results excluded from an incompatible interpretation.}
\label{fig:evidence-status}
\end{figure}
